% Part: computability
% Chapter: recursive-functions
% Section: pr-relations

\documentclass[../../../include/open-logic-section]{subfiles}

\begin{document}

\olfileid{cmp}{rec}{prr}
\olsection{العلائق العودية البدائية}

\begin{defn}
العلاقة $R(\vec x)$ عودية بدائية متى كانت دالتها المميزة
\[
\Char{R}(\vec x) = \left\{
  \begin{array}{ll}
  1 & \mbox{إذا صدقت $R(\vec x)$} \\
  0 & \mbox{في غير ذلك}
  \end{array}
\right.
\]
عودية بدائية.
\end{defn}

فمعنى القول بعودية العلاقة $R(\vec x)$ البدائية أنها
تأخذ صورة $\Char{R}(\vec x)=1$، وأن $\Char{R}$ دالة
عودية بدائية لا تعطي، لأي دخل، إلا $1$ أو $0$. ومن أمثلة ذلك العلاقة
$\fn{IsZero}(x)$، الصادقة إذا وفقط إذا كان $x=0$؛ فإن دالتها المميزة هي
$\Char{\fn{IsZero}}$، وتعريفها بالعودية البدائية كما يأتي:
\begin{align*}
\Char{\fn{IsZero}}(0) & = 1,\\
\Char{\fn{IsZero}}(x+1) & = 0.
\end{align*}

ويجوز تركيب العلائق مع دوال عودية بدائية
أخرى، فلا تخرج بذلك عن العودية البدائية. ومن هذا نحصل على العلاقتين الآتيتين:
\begin{enumerate}
\item علاقة المساواة، $x=y$، المعرفة بـ
$\fn{IsZero}(\left|x-y\right|)$.
\item علاقة الأصغر من أو المساوي، $x\leq y$، المعرفة بـ
$\fn{IsZero}(x\tsub y)$.
\end{enumerate}

\begin{prop}
  تحفظ العمليات البوليانية العودية البدائية للعلائق؛ فمتى
  كانت $P(\vec x)$ و$Q(\vec x)$ عوديتين بدائيتين، لزمت العودية البدائية
  للعلائق الآتية:
  \begin{enumerate}
  \item $\lnot P(\vec x)$
  \item $P(\vec x) \land Q(\vec x)$
  \item $P(\vec x) \lor Q(\vec x)$
  \item $P(\vec x) \lif Q(\vec x)$
  \end{enumerate}
\end{prop}

\begin{proof}
  لنفرض عودية $P(\vec x)$ و$Q(\vec x)$ البدائية؛ فدالتاهما
  المميزتان $\Char{P}$ و$\Char{Q}$ عوديتان بدائيتان. والمطلوب إثبات العودية
  البدائية للدوال المميزة لـ$\lnot P(\vec x)$ ولسائر العلائق المذكورة.
  \[
  \Char{\lnot P}(\vec x) = \begin{cases}
    0 & \text{إذا كانت $\Char{P}(\vec x) = 1$}\\
    1 & \text{في غير ذلك}
  \end{cases}
  \]
  ويكفي لتعريف $\Char{\lnot P}(\vec x)$ أن نجعلها
  $1\tsub\Char{P}(\vec x)$.
  \[
  \Char{P \land Q}(\vec x) = \begin{cases}
    1 & \text{إذا كانت $\Char{P}(\vec x) = \Char{Q}(\vec x) = 1$}\\
    0 & \text{في غير ذلك}
  \end{cases}
  \]
  وأما $\Char{P\land Q}(\vec x)$ فتعريفها أن نجعلها
  $\Char{P}(\vec x)\cdot\Char{Q}(\vec x)$، أو بأنها
  $\fn{min}(\Char{P}(\vec x),\Char{Q}(\vec x))$. وبالمثل،
  \begin{align*}
    \Char{P \lor Q}(\vec x) & = \fn{max}(\Char{P}(\vec x), \Char{Q}(\vec x)) \text{، و}\\
    \Char{P \lif Q}(\vec x) & = \fn{max}(1 \tsub
  \Char{P}(\vec x), \Char{Q}(\vec x)).
  \end{align*}
\end{proof}

\begin{prop}
  والتكميم المحدود كذلك يحفظ العودية البدائية؛ فإذا
  كانت العلاقة $R(\vec x,z)$ عودية بدائية، كانت العلاقتان
  \begin{align*}
    & \bforall{z < y}{R(\vec x, z)} \text{، و}\\
    & \bexists{z < y}{R(\vec x, z)}
  \end{align*}
  عوديتين بدائيتين. وتصدق $\bforall{z<y}{R(\vec x,z)}$ على $\vec x$
  و$y$ إذا وفقط إذا صدقت $R(\vec x,z)$ لكل~$z$ أصغر من~$y$، وكذلك
  الحال في $\bexists{z<y}{R(\vec x,z)}$.
\end{prop}

\begin{proof}
  اصطلاحنا أن $\bforall{z<0}{R(\vec x,z)}$ صادقة، إذ ليس
  ثمة $z$ أصغر من~$0$، وأن
  $\bexists{z<0}{R(\vec x,z)}$ كاذبة. والمكمم الكلي المحدود بمنزلة
  جداء منته أو تكرار لأخذ الأصغر؛ فمتى وضعنا
  $P(\vec x,y)\defiff\bforall{z<y}{R(\vec x,z)}$، عرّفنا
  $\Char{P}(\vec x,y)$ هكذا:
  \begin{align*}
    \Char{P}(\vec x, 0) & = 1\\
    \Char{P}(\vec x, y+1) & =
    \fn{min}(\Char{P}(\vec x, y), \Char{R}(\vec x, y)).
  \end{align*}
  وأما التكميم الوجودي المحدود فيعرّف بمثل ذلك باستعمال $\fn{max}$. ولنا أيضًا
  أن نرده إلى التكميم الكلي المحدود، اعتمادًا على التكافؤ
  $\bexists{z<y}{R(\vec x,z)}\liff
  \lnot\bforall{z<y}{\lnot R(\vec x,z)}$. ومن المفيد هنا أن المكمم
  المحدود الذي صورته $\bexists{x\leq y}{\dots x\dots}$ يكافئ
  $\bexists{x<y+1}{\dots x\dots}$.
\end{proof}

\begin{prob}
  أثبت العودية البدائية للعلاقة الثلاثية الحجج $x\equiv y\mod n$ (التطابق بترديد
  ~ $n$).
\end{prob}

ومما ينفعنا من الدوال العودية البدائية الدالة الشرطية
$\fn{cond}(x,y,z)$ المعرفة بـ
\begin{align*}
  \fn{cond}(x,y,z) & = \begin{cases}
  y & \text{إذا كان $x=0$} \\
  z & \text{في غير ذلك}.
\end{cases}
\intertext{وهي معرفة عوديًا بـ}
\fn{cond}(0,y,z) & = y,\\
\fn{cond}(x+1,y,z) & = z.
\end{align*}
وبها يتبين جواز تعريف الدوال العودية البدائية بالحالات، متى كانت
العلائق التي تفصل تلك الحالات عودية بدائية:

\begin{prop}
إذا كانت $g_0(\vec x)$، …،~$g_m(\vec x)$ دوال عودية بدائية، وكانت
$R_0(\vec x)$، …،~$R_{m-1}(\vec x)$ علائق عودية بدائية، فإن الدالة
$f$ المعرفة بـ
\[
f(\vec x) = \begin{cases}
    g_0(\vec x) & \text{إذا صدقت $R_0(\vec{x})$} \\
    g_1(\vec x) & \text{إذا صدقت $R_1(\vec{x})$ ولم تصدق $R_0(\vec{x})$} \\
    \vdots & \\
    g_{m-1}(\vec x) & \text{إذا صدقت $R_{m-1}(\vec{x})$ ولم تصدق أي
      من سابقاتها}
    \\
    g_m(\vec x) & \mbox{في غير ذلك}
\end{cases}
\]
عودية بدائية أيضًا.
\end{prop}

\begin{proof}
  أما عند $m=1$ فالدالة المطلوبة عين المعرّفة بـ
  \[
  f(\vec x) = \fn{cond}(\Char{\lnot R_0}(\vec x),g_0(\vec x),g_1(\vec
  x)).
  \]
  وأما إذا زاد $m$ على~$1$، ركبنا تعريفات من هذا القبيل فحصل المطلوب.
\end{proof}
\end{document}
